\section{Introduction}
\label{sec:intro}

\IEEEPARstart{A}{lert} triage is the bottleneck of security operations. Analysts in security operations
centres describe being flooded with alarms, most of which are false positives, and spending much of their
time confirming that activity was benign~\cite{alahmadi2022falsepositives,kokulu2019soc}. The resulting
workload is a recognised cause of analyst burnout~\cite{sundaramurthy2015burnout}. Detection research has
responded by building systems that reason over \emph{provenance graphs}, the causal record of processes,
files, sockets and their interactions collected by endpoint sensors, to reconstruct attacks and to separate
them from benign activity~\cite{hossain2017sleuth,milajerdi2019holmes,hassan2019nodoze,han2020unicorn,cheng2024kairos}.
These systems raise or rank alerts, but the investigation that follows, deciding what an alert means and
what must be done about it, is still largely manual.

Large language models (LLMs) used as agents, interleaving reasoning with tool calls~\cite{yao2023react,schick2023toolformer},
are a natural fit for that investigation. An agent that receives an alert can query the evidence graph,
follow relations from the triggering process to its parent, children, written files and network
connections, and submit a verdict with the entities that require a response. Recent industrial evaluations
frame exactly this task: the BlueSec1 competition runtime~\cite{bluesec1_reference} gives an agent an alert
and a small tool interface over an evidence graph (\tool{search}, \tool{get\_entity}, \tool{get\_relation},
\tool{finish\_investigation}) and scores both the quality of the submission and the number of tool calls
spent.

Building such an agent raises problems that general agent benchmarks~\cite{liu2024agentbench,jimenez2024swebench}
and LLM security benchmarks~\cite{bhatt2023cyberseceval,alam2024ctibench,zhang2025cybench} do not address.
First, \emph{cost}: in a SOC, and in budget-scored evaluations, every lookup counts, so an agent must decide
which evidence is worth fetching. Second, \emph{grounding}: an LLM may cite entities that never appeared in
any tool result, which is unacceptable for incident response. Third, \emph{robustness}: alerts that look the
same can mean opposite things (an operator stopping a web service versus a package manager starting it), so
an agent must clear benign activity as reliably as it confirms attacks. Fourth, \emph{evaluation}: public
provenance datasets either lack entity-level ground truth or are hundreds of gigabytes in size, and none
packages alert-investigation tasks with benign look-alikes across operating systems. Fifth,
\emph{shortcuts}: public attack recordings are labelled by how they were made, and the labels leak into dataset
names, lab host names and the names operators chose for their tools. An agent can then be right for the wrong
reason, and an intruder who avoids such names would not be caught the same way, the evasion problem known from
mimicry attacks on host-based detection~\cite{wagner2002mimicry,goyal2023mimicry}.

This paper makes four contributions.

\begin{enumerate}
\item \textbf{A grounded, budget-aware investigation agent} (\S\ref{sec:method}). Tools are derived from the
runtime's advertised catalog. Every call passes through local safeguards that answer repeated, malformed,
ungrounded or over-budget calls without spending a runtime call: a response cache, schema validation, a
grounding check on the submission, and soft and hard budgets with deadline control. Tasks run in parallel
under a runtime limit the agent learns from refusals. We also report a practical failure mode, model-side
refusals of requests that ask the model to externalise its reasoning, and a field-aliasing fix that keeps the
runtime contract intact.
\item \textbf{\bench, a cross-platform benchmark from public telemetry} (\S\ref{sec:benchmark}). A graph
builder converts Windows Sysmon and Sysmon-for-Linux recordings from OTRF Security-Datasets~\cite{otrf_securitydatasets}
and Splunk attack\_data~\cite{splunk_attackdata} into 21 evidence-graph tasks: 15 ATT\&CK technique
executions and 6 benign look-alikes from the same hosts, with ground truth defined as selectors over the
graph and a scorer for verdict, response artifacts, legitimacy evidence and efficiency. The benchmark is a
0.55\,MB file and can be rebuilt from the original recordings. Two cue-controlled variants share its graphs
and answers: \emph{blind} removes label cues of the recordings, and \emph{mimicry} also renames
attacker-chosen artefacts without changing behaviour.
\item \textbf{An empirical evaluation and ablation protocol} (\S\ref{sec:experiments}--\S\ref{sec:results}).
In a pilot on the \ResPilotTasks{} Windows tasks with Claude Opus 5.5, the agent reached the correct verdict on
all of them, including all three benign tasks, with mean quality \ResPilotQuality{} and mean efficiency
\ResPilotEfficiency{}. The ablation protocol switches off one safeguard at a time and reports quality,
efficiency, verdict accuracy, false-positive and false-negative rates, and cost, per platform and per verdict
class, with paired task-level confidence intervals.
\ifnum\ResStudyAvailable=1
On all 21 tasks the full agent reached reward \ResFullReward{} with no false positives. The
investigation-method prompt is what makes it efficient (\ResNoMethodCalls{} instead of \ResFullCalls{} calls
per task without it);
the safeguards did not change reward measurably with this model, a negative result we report with its
statistics. Per-call traces explain it and expose the actual sources of error: identifiers of the wrong type
that a substring grounding check accepts, an efficiency instruction that withheld a needed identifier, and label
cues cited in \TrAllSourceCuePct\% of submissions.
\fi
\item \textbf{An open, reproducible release.} Agent, benchmark, traces of every tool call, and the script that
regenerates all tables and plots from run folders.
\end{enumerate}

The remainder of the paper reviews related work (\S\ref{sec:background}), defines the task
(\S\ref{sec:problem}), presents the agent (\S\ref{sec:method}) and the benchmark (\S\ref{sec:benchmark}),
describes the experimental setup (\S\ref{sec:experiments}), reports results (\S\ref{sec:results}), and
discusses implications (\S\ref{sec:discussion}) and limitations (\S\ref{sec:limitations}).
